\printedpage{130}
% head: A. Grothendieck : Dérivateurs, Ch. XII, Édition des Universités de Montpellier II et Paris VII

\[
\begin{tikzcd}[column sep=small, row sep=small]
 & C(X') & \\
C(X) \arrow[ur, "C(i')"] & & X' \arrow[ul, "i_{X'}"'] \\
 & X \arrow[ul, "i_X"] \arrow[ur, "i'"'] &
\end{tikzcd}
\]

is commutative (by functoriality of $i_X : X \longrightarrow C(X)$), and we apply this to the case of $X' = C'(X)$, $i' = i'_X$.

This remark encourages us to make a leap into abstraction, and to introduce the category\msp{page 198}
\[
\mathcal{C} = \mathcal{C}_S = \begin{cases}
\text{category of pairs } (C,i), \text{ with} \\
C : \Cat/S \longrightarrow \Cat/S \text{ factoring through } \Fib_W S, \\
i : \id_{\Cat/S} \longrightarrow C \text{ such that } i_X \in W \ \forall\ X \text{ in } \Cat/S,
\end{cases}
\tag*{$(*)$}
\]
an arrow
\[
(C,i) \longrightarrow (C',i')
\]
being an arrow
\[
C \xrightarrow{u} C' \qquad \text{(in } \Hom(\Cat/S, \Cat/S))
\]
making the square \texttt{[read: the diagram]} commutative
\[
\begin{tikzcd}[column sep=small]
 & \id_{\Cat/S} \arrow[dl, "i"'] \arrow[dr, "i'"] & \\
C \arrow[rr, "u"] & & C'.
\end{tikzcd}
\]

In other words, $\mathcal{C}$ is defined as a full subcategory of % [sic] sous-catégiorie -> sous-catégorie, misspelled word
\[
\id_{\Cat/S}\backslash\Hom(\Cat/S, \Cat/S),
\]
namely the one whose objects $(C,i)$ satisfy the two conditions $(*)$ above (both involving the localizer $W$).

\msp{page 199}These conditions imply this

\begin{enumerate}
\item[a)] $\forall\ X$ in Cat/$S$, $(u_X : C(X) \longrightarrow C'(X)) \in W$, since
\[
\begin{tikzcd}[column sep=small]
 & X \arrow[dl, "i_X\in W"'] \arrow[dr, "i'_X\in W"] & \\
C(X) \arrow[rr, "u_X"] & & C'(X),
\end{tikzcd}
\]
\printedpage{131}
% head: A. Grothendieck : Dérivateurs, Ch. XII, Édition des Universités de Montpellier II et Paris VII
and

\item[b)] if $W$ satisfies W(4b), so that $u_X \in W^{\mathrm{f}}_S$, then this defines an \textit{isomorphism} between $\varphi_C(X)$ and $\varphi_{C'}(X)$ % [sic] ce définit -> ceci définit, missing word
\[
\begin{tikzcd}[row sep=small]
s \arrow[r, mapsto] & \hot_W(C(X)_s) \arrow[d, "\wr\, u_{X,s}"] \\
s \arrow[r, mapsto] & \hot_W(C'(X)_s).
\end{tikzcd}
\]
\end{enumerate}

Better still, in this case, the composite functor
\[
\mathcal{C} \times \Cat/S \longrightarrow \Fib_W S \longrightarrow \HOT^{\mathrm{lc}}_W(S)
\]
gives rise to
\[
\mathcal{C} \longrightarrow \Hom(\Cat/S, \HOT^{\mathrm{lc}}_W(S))
\]
which transforms every arrow of $\mathcal{C}$ into \texttt{[an]} \textit{isomorphism,} hence factors through the fundamental groupoid of $\mathcal{C}$,
\[
\Pi\mathcal{C} \longrightarrow \Hom(\Cat/S, \HOT^{\mathrm{lc}}_W(S)).
\]

\msp{page 200}As for the functor considered earlier
\[
\begin{aligned}
\mathcal{C} &\quad\longrightarrow\quad \underbrace{\Hom\,(\Cat/S, \Hom(S^\circ, \Hot_W))}_{\simeq\Hom((\Cat/S)\times S^\circ, \Hot_W)} \\
C &\quad\longmapsto\quad ((X,s) \longmapsto \hot_W(C(X)_s))\,,
\end{aligned}
\]
it is deduced from the preceding one by composing with
\[
\begin{tikzcd}[column sep=small]
\HOT_W(S) \arrow[rr] \arrow[dr, dotted] & & \Hom(S^\circ, \Hot_W), \\
 & \Hom(\Pi S^\circ, \Hot_W) \arrow[ur, dotted] &
\end{tikzcd}
\]
more exactly, with the functor deduced from it
\[
\Hom(\Cat/S, \HOT_W(S)) \longrightarrow \Hom(\Cat/S, \Hom(S^\circ, \Hot_W)).
\]

We see emerging here a plausible principle of canonicity of isomorphisms, as regards the functors
\[
C_W : \Cat/S \longrightarrow \HOT^{\mathrm{lc}}_W(S)
\]
defined by the $C \in \mathrm{Ob}\,\mathcal{C}$, by the relation
\[
\Pi\mathcal{C} \approx \ \text{one-point category, i.e. } \mathcal{C} \text{ is 1-connected.}
\]

Moreover, the arguments used so far are of so formal a nature, that the specific aspects of the Modelizer Cat
\printedpage{132}
% head: A. Grothendieck : Dérivateurs, Ch. XII, Édition des Universités de Montpellier II et Paris VII
\msp{page 201}seem to play practically no role in them, or a very unobtrusive role.

The square on page 197 shows above all that the category $\mathcal{C}$ is filtering - $(C,i)$, $(C,i')$ both map into $(C'',i'')$, with $C'' = C \circ C'$, and $i'' = (C * i') \circ i = (i * C') \circ i'$. There is even a law of composition in $\mathcal{C}$, by no means commutative or unital, that's for sure, but surely associative, and (if this composition is denoted by $*$) with functorial morphisms % [sic] (C,i') -> (C',i'), the second couple of the following formula is (C',i')
\[
\xi \xrightarrow{\alpha_{\xi,\xi'}} \xi * \xi' \xleftarrow{\beta_{\xi,\xi'}} \xi'
\]
which ensure the upward filtering character of $\mathcal{C}$. These arrows seem somehow to replace the missing two-sided unit, and one expects them to have the same excellent properties as if one had such a unit, concerning the double composites $(\xi * \xi') * \xi'' \simeq \xi * (\xi' * \xi'')$,

\msp{page 202}
\[
(\xi * \xi') * \xi'' \xrightarrow{\sim} \xi * (\xi' * \xi'').
\]

There are two ways of mapping a binary product
\[
\xi * \xi', \qquad \xi' * \xi'', \qquad \xi * \xi''
\]
into the ternary product, and well, they must in each of the three cases, be equal. And likewise there are, this granted, again two ways of mapping $\xi$ into the ternary product, factoring either through $\xi * \xi'$, or through $\xi * \xi''$, and likewise for the other two functors $\xi'$, $\xi''$; and there again, these two ways must be equal, in each of these three cases - so $3 + 3 = 6$ equalities in all to write down and to check - we've seen worse!

The computational philosophy should be this: If we have an $n$-fold product
\[
\xi_1 * \xi_2 * \cdots * \xi_n
\]
(where I dispense with putting in parentheses \dots), for every partial $m$-fold product $(m \in [1, n-1])$\msp{page 203}
\[
\xi_{i_1} * \xi_{i_2} * \cdots * \xi_{i_m} \qquad\qquad (i_1 < i_2 < \cdots < i_m),
\]
using the associativity isomorphisms and the $\alpha_{\xi,\xi'}$, $\beta_{\xi,\xi'}$ left and right, there is nevertheless only one way of mapping it into the $n$-fold product. I really have the impression that this must be true. (Does this already imply that $\mathcal{C}$ is 1-connected?)

In fact, there is indeed a two-sided unit, but it is not in $\mathcal{C}$, but in the larger category
\[
\id_{\Cat/S}\backslash\Hom(\Cat/S, \Cat/S),
\]
which itself, like any category
\[
\id_{\mathcal{X}}\backslash\Hom(\mathcal{X}, \mathcal{X}),
\]
\printedpage{133}
% head: A. Grothendieck : Dérivateurs, Ch. XII, Édition des Universités de Montpellier II et Paris VII
is a monoidal category, i.e. with an associative \textit{and} unital product. The arrows $\alpha_{\xi,\xi'}$, $\beta_{\xi,\xi'}$ are the arrows that result from the presence of this unit (seeing that $\mathcal{C}$ is a full subcategory), and then it should become clear that we have the said compatibilities.

\msp{page 204}But even a monoidal category, is it 1-connected? For example $B_G$, for $G$ a group say, is indeed a monoidal category and even ``groupal'', a $G$-category (as strict as can be), and yet it is \textit{not} 1-connected!

Of course, we can short-circuit the question of the 1-connectedness of $\mathcal{C}$, by noting that after all we already have a priori a privileged object in
\[
\Hom(\Cat/S, \HOT^{\mathrm{lc}}_W(S)),
\]
i.e. a canonical functor
\[
\Cat/S \longrightarrow \HOT^{\mathrm{lc}}_W(S),
\]
namely the composite functor
\[
\Cat/S \underbrace{\longrightarrow}_{\substack{\text{localization}\\ \text{by } W_S}} (\Cat/S)W_S^{-1} \longrightarrow \underbrace{\HOT^{\mathrm{lc}}_W(S)}_{\overset{\text{def}}{=} (\Fib_W S)(W^{\mathrm{f}}_S)^{-1}} \quad ,
\tag*{$(*)$}
\]
where the second arrow is defined as the quasi-inverse of
\[
(\Fib_W S)(W^{\mathrm{f}}_S)^{-1} \longrightarrow (\Cat/S)W_S^{-1}
\tag*{$(**)$}
\]

\msp{page 205}deduced from the inclusion
\[
(\Fib_W S, W^{\mathrm{f}}_S) \hookrightarrow (\Cat/S, W_S).
\]

The fact that $(*)$ is an equivalence, is proved precisely by the fact of the existence of a functor $C : \Cat/S \longrightarrow \Cat/S$ and of an $i : \id_{\Cat/S} \longrightarrow C$, having the properties we know, i.e. from the fact that $C \neq \emptyset$. Moreover, the datum of an object $C$ of $\mathcal{C}$ gives a convenient way of making the quasi-inverse of the said functor explicit. But this quasi-inverse has an existence that is canonical and independent of the choice of a $C$, i.e. of a mode of calculation. The fact that the composite functor $\varphi$ in $(*)$ can be ``lifted'' (up to isomorphism) to a functor % [sic] C \neq \emptyset -> \mathcal{C} \neq \emptyset, the non-emptiness concerns the category \mathcal{C}
\[
C : \Cat/S \longrightarrow \Fib_W S
\]
between $\varphi(X) = \hot_{W,S}(X)$ and $\hot_{W,S}(C(X))$ and that the isomorphism can be made explicit, it too, by an $X$-functorial arrow
\[
i_X : X \longrightarrow C(X), \qquad i_X \in W,
\]
in Cat/$S$ itself (and not only in the localized $(\Cat/S)W_S^{-1}$),
\printedpage{134}
% head: A. Grothendieck : Dérivateurs, Ch. XII, Édition des Universités de Montpellier II et Paris VII
\msp{page 206}is a particular virtue of the \texttt{[illegible word]}, which should not for all that make us forget the intrinsic character of the arrows $(*)$, independently of any mode of calculation.

{\reversemarginpar\marginnote{$\|$}}% double vertical bar printed in the left margin beside the first lines of this paragraph
The key to the situation, and to the questions of canonicity, seems to me to be in the fact that $(**)$ is an equivalence of categories - and one must admit that, however familiar and almost banal it has finally become to me, this is a quite remarkable fact and by no means obvious a priori. (But probably one can also get it by the general \textit{fourbis} (gear) à la \textsc{Quillen-Thomason}, without fighting with a formalism of paths \dots) The ``canonical isomorphisms'' that we can establish between a ``theory $C$ '' and a ``theory $C'$ '' (via two different constructions, e.g., of categories of paths, or via the \textsc{Quillen-Thomason} approach, without paths) must always come back to those

\msp{page 207}that we can define via the comparison of the ``theory $C$ '' on the one hand, the ``theory $C'$ '' on the other, with what we get directly from the composite arrow $(*)$ (p. 204). Finally, it seems to me that \textit{the existence} of a theory of the ``homotopy fibers'' (or ``$W$-homotopic'') of an arrow $X \longrightarrow S$ in Cat, results from the equivalence $(**)$ (p. 204), and this theory could probably to a certain extent be developed from there, without a choice of a functor $C$ or of a model category structure in the sense of \textsc{Quillen}. % [sic] qu'on peu définir -> qu'on peut définir, verb form

\section*{4\quad Axiomatization of ``$W$-homotopic fibers''.}

Yesterday's reflection makes me return to the axiomatics in ``categories of models'' % [sic] reflexion -> réflexion, missing accent
\[
(\mathcal{M}, W)
\]
(\footnote[119]{$W$ satisfies the axiom of weak saturation: isomorphisms \texttt{[are]} in $W$, and if two among $f$, $g$, $gf$ \texttt{[are]} in $W$, the third also.}) more general than Cat, in the spirit of VI (cf. theorem-scholie 2, p. 108-113). For the moment, we are going to suppose only $\mathcal{M}$ stable under fiber products (we do not need a final object). We introduce the set
\[
\mathrm{Fib}_W \subseteq \mathrm{Fl}(\mathcal{M})
\]

\msp{page 208}of \textit{$W$-fibrations} (``ultrastrong'' being understood!) by the condition

\begin{quote}
$f : X \longrightarrow Y$ is in $\mathrm{Fib}_W$ if and only if the base change functor by $X$,
\[
\begin{aligned}
Y' &\quad\longmapsto\quad X' = X \times_Y Y' \\
\mathcal{M}/Y &\quad\longrightarrow\quad \mathcal{M}/X
\end{aligned}
\]
\end{quote}